\section{Quadratic mass growth}\label{sec:mass}

Let \(P,\mu\) be as in Proposition~\ref{prop:separation}.
We use the annihilation of closed invariant forms to bound the
mass of \(P\) in every small ball.
Write \(B_s(x)\) for a geodesic ball.
\begin{proposition}\label{prop:growth}
There is a constant \(C\) such that
\begin{equation}\label{eq:growth}
\mu(B_s(x))\le Cs^2
\end{equation}
for all \(x\in X\) and \(s>0\).
\end{proposition}

The proof uses closed invariant test forms
\[
Du=(1-BR)dd_J^cu,\qquad d_J^cu=-du\circ J.
\]
The Hessian terms of \(dd_J^cu\), with \(J\) frozen at a point,
form an invariant two-form. Consequently \(Rdd_J^c\) is a
first-order differential operator on functions.

We first spell out the local estimate needed for the
order-zero correction. In dimension four, an order-zero
pseudodifferential operator has an off-diagonal kernel bounded
by \(C|y-z|^{-4}\) in local coordinates
\cite[Section~0.2, Proposition~0.2.B]{Taylor1991}.
Work in a fixed smooth local bundle trivialization.
If \(f\) is supported in a coordinate ball of radius \(a\le1\)
and its coefficient derivatives satisfy
\begin{equation}\label{eq:scaled-input}
|\partial^jf|\le C_jMa^{-j},
\end{equation}
then its local contribution to the operator is bounded by
\(CM\). Indeed, write \(f(z)=M\widetilde f((z-z_0)/a)\).
Six bounded derivatives of \(\widetilde f\) give
\(|\widehat{\widetilde f}(\xi)|\le C(1+|\xi|^2)^{-3}\)
by integration by parts. This bound is integrable in four
dimensions, so \(\|\widehat{\widetilde f}\|_1\le C\), and hence
\(\|\widehat f\|_1\le CM\). The bounded-symbol Fourier formula
gives the assertion; smoothing remainders are controlled by
\(\|f\|_1\). This estimate uses the derivative bounds
\eqref{eq:scaled-input}, not general \(L^\infty\) boundedness.

\begin{lemma}\label{lem:radial-correction}
Choose normal coordinates centered at \(x\), on a uniformly
small fixed ball, and put \(t=d_g(x,y)\). Let a smooth radial
cutoff be one on a smaller ball and vanish outside the
coordinate ball. For \(s>0\) small, let \(u_s,h_s\) be this
cutoff times
\[
\log\sqrt{s^2+t^2},\qquad \sqrt{s^2+t^2},
\]
respectively, extended by zero. Near the center,
\begin{equation}\label{eq:correction-bounds}
|BRdd_J^cu_s|\le \frac{C}{s+t},
\qquad
|BRdd_J^ch_s|\le C(1+|\log(s+t)|).
\end{equation}
Both expressions are uniformly bounded at any fixed positive
distance from \(x\).
\end{lemma}

\begin{proof}
The first-order property of \(Rdd_J^c\) gives
\[
|\partial^j(Rdd_J^cu_s)|\le C_j(s+t)^{-1-j},\qquad
|\partial^j(Rdd_J^ch_s)|\le C_j(s+t)^{-j}
\]
near \(x\). Away from \(x\), all fixed-order derivatives are
uniformly bounded, including terms from the cutoff.

At \(y\), put \(a=s+t\), assuming \(a\) small. Localize the
input to a ball about \(y\) of radius a small fixed multiple
of \(a\). On this ball \(s+d_g(x,\cdot)\) is comparable to
\(a\), so \eqref{eq:scaled-input} gives contributions \(C/a\)
and \(C\). For the rest of the input within distance \(O(a)\)
of \(x\), the kernel is bounded by \(Ca^{-4}\). The respective
input \(L^1\) bounds on this region are \(Ca^3\) and \(Ca^4\),
again giving \(C/a\) and \(C\).

On the more distant shells about \(x\), with radius
\(\rho\ge Ca\), the kernel is \(O(\rho^{-4})\), the input
bounds are \(C/\rho\) and \(C\), and the shell volume is
\(O(\rho^4)\). Summation over dyadic shells gives \(C/a\)
and \(C(1+|\log a|)\). The logarithmic input has uniformly
bounded total \(L^1\) norm, since its local integral is at most
\(C\int_0^{\rho_0}\rho^3/(s+\rho)\dd\rho\le C\rho_0^3\).
The square-root input has bounded \(L^1\) norm as well.
Thus the remaining smooth kernel terms are bounded uniformly.
The same decomposition gives uniform bounds
when \(y\) stays a fixed distance from the center.
Compactness makes all chart and operator constants uniform
in \(x\).
\end{proof}

\begin{proof}[Proof of Proposition~\ref{prop:growth}]
In the normal coordinates of Lemma~\ref{lem:radial-correction},
let \(J_0=J(x)\) be the constant complex structure, orthogonal
for the Euclidean metric at the center. On the inner ball,
the Hessians of the two radial functions are
\[
\frac{I}{s^2+t^2}-\frac{2zz^\top}{(s^2+t^2)^2},
\qquad
\frac{I}{\sqrt{s^2+t^2}}-\frac{zz^\top}{(s^2+t^2)^{3/2}}.
\]
Since
\((z\cdot v)^2+(z\cdot J_0v)^2\le t^2|v|^2\), their
complex-line traces satisfy
\begin{align}
dd^c_{J_0}u_s(v,J_0v)
&\ge\frac{2s^2|v|^2}{(s^2+t^2)^2},\label{eq:log-trace}\\
dd^c_{J_0}h_s(v,J_0v)
&\ge\frac{|v|^2}{\sqrt{s^2+t^2}}.\label{eq:sqrt-trace}
\end{align}
The first lower bound is nonnegative everywhere and at least
\(c/s^2\) for \(t<s\) and \(g\)-unit \(v\).

Replacing \(J_0\) by \(J(y)\) both in the differential expression
and in the second argument costs at most \(C/(s+t)\) for
\(u_s\) and \(C\) for \(h_s\). To see this, use
\(J(y)-J_0=O(t)\) on the Hessian terms and the bounded first
derivatives of \(J\) on the gradient terms. Euclidean and
\(g\)-norms are uniformly comparable.

Together with \eqref{eq:correction-bounds}, this proves,
for \(s+t\) sufficiently small,
\[
Du_s(v,Jv)\ge \frac{c}{s^2}\ind_{\{t<s\}}-\frac{C}{s+t},
\qquad
Dh_s(v,Jv)\ge \frac{c'}{s+t}
\]
with \(c,c'>0\). In the second estimate the logarithmic
correction is absorbed by the positive inverse-distance
term on a sufficiently small fixed neighborhood.
Choose a fixed \(A>0\) large enough to absorb the negative
term in the first estimate. Outside that neighborhood all
terms are uniformly bounded. We obtain globally
\[
D(u_s+Ah_s)(v,Jv)\ge
\frac{c}{s^2}\ind_{B_s(x)}-C'
\]
for all unit \(v\) and small \(s\).

The form on the left is closed and invariant. Applying \(P\),
which annihilates it, and using \(\mu(X)=1\), gives
\(0\ge cs^{-2}\mu(B_s(x))-C'\).
This proves \eqref{eq:growth} for small \(s\). Increasing
the constant handles all larger radii.
\end{proof}
